%In the beginning of this section we highlighted the fact that the $(2n)^{\mathrm{th}}$ slice of $\MU_\R$ is a sum of copies of $\Sigma^{n \rho} \uZ$, which is \emph{not} in (a shift of) the heart of the usual $t$-structure. From this we were able to conclude that the obvious comparison functor between $C\ta$-modules and graded equivariant objects was inadiqute for our purposes.
%In this subsection, we discuss twisted $t$-structures on filtered and graded categories. Using this, we will produce ``twisting functors'' which will allow us to . We will take particular care to study how these twisting functors interact with symmetric monoidal structures.
%
We begin this subsection with a discussion of several natural $t$-structures on categories of graded and filtered objects, obtained by twisting a given $t$-structure by a Picard element.
In the filtered case our constructions are a straightforward generalization of a construction of Beilinson. Although these constructions are simple, they have the effect of modifying the symmetric monoidal structure on the heart of a category.
Next, we explicitly describe this modification at the level of the heart, where it admits a description in terms of an Euler characteristic. Read another way, there is an Euler characteristic obstruction to producing symmetric monoidal twisting functors. 
We close the subsection by producing symmetric monoidal twisting functors whenever this Euler characteristic obstruction vanishes.


\begin{dfn} \label{dfn:L-twist}
  For the remainder of this subsection, we fix the following data:
  a stable, presentably symmetric monoidal category $\mathcal{C}$
  equipped with a compatible $t$-structure\footnote{Here compatible means that a tensor product of objects which are $\geq 0$ is $\geq 0$ itself, and that the unit is $\geq 0$.}
  and a Picard element $\mathcal{L}$.

  We then define the $\mathcal{L}$-twisted $t$-structure on $\mathcal{C}^{\Gr}$ and $\mathcal{C}^{\Fil}$ as follows:\todo{Cite something to show this allowed?}
  \begin{itemize}
  \item An object $X_\bullet$ of $\mathcal{C}^{\Gr}$ is $\geq 0$ in the $\mathcal{L}$-twisted $t$-structure if, for all $n$, $\mathcal{L}^{\otimes -n} \otimes X_n$ is $\geq 0$ in the original $t$-structure.
  \item Similarly, an object $X_\bullet$ of $\mathcal{C}^{\Fil}$ is $\geq 0$ in the $\mathcal{L}$-twisted $t$-structure if $\mathcal{L} ^{\otimes -n} \otimes X_n$ is $\geq 0$ in the original $t$-structure, for all $n$.
  \end{itemize}
  We denote the hearts of these $t$-structures by $\CC^{\Gr, \mathcal{L}, \heartsuit}$ and $\CC^{\Fil, \mathcal{L}, \heartsuit}$ respectively.
  We denote the $i^{\mathrm{th}}$ $\mathcal{L}$-twisted $t$-structure homotopy objects by $\piu_{i} ^{\mathcal{L}, \Gr}$ and $\piu_{i} ^{\mathcal{L}, \Fil}$ respectively.
  When $\mathcal{L} = \o$, we frequently omit it from the notation.
\end{dfn}

We summarize the basic properties of this definition in the following lemma.

\begin{lem} \label{prop:twist-grade} \label{prop:twist-filt} \label{prop:twist-1-conn}
  In the graded case:
  \begin{enumerate}
  \item The $\mathcal{L}$-twisted $t$-structure on $\CC^{\Gr}$ is compatible with the symmetric monoidal structure.
  \item An object $X_\bullet$ is $< 0$ in the $\mathcal{L}$-twisted $t$-structure if and only if, for all $n$, $\mathcal{L}^{\otimes -n} \otimes X_n$ is $< 0$ in the original $t$-structure.
  \item The $t$-structure homotopy groups are determined by the following formula:
\[ (\underline{\pi}_0^{\mathcal{L}, \Gr} X_\bullet )_n \cong \mathcal{L}^{\otimes n} \underline{\pi}_0(\mathcal{L}^{\otimes -n} \otimes X_n). \]
  \end{enumerate}
  Assuming that $\mathcal{L} \geq 0$, we obtain similar results in the filtered case:
  \begin{enumerate}
  \item[(1')] The $\mathcal{L}$-twisted $t$-structure on $\CC^{\Fil}$ is compatible with the symmetric monoidal structure.
  \item[(2')] An object $X_\bullet$ is $< 0$ in the $\mathcal{L}$-twisted $t$-structure if and only if $\mathcal{L}^{\otimes -n} \otimes X_n < 0$.
  \item[(3')] The $t$-structure homotopy groups are given the by following formula:
  \[ (\underline{\pi}_0^{\mathcal{L}, \Fil} X_\bullet )_n \cong \mathcal{L}^{\otimes n} \underline{\pi}_0(\mathcal{L}^{\otimes -n} \otimes X_n). \]

  \end{enumerate}

  Under the stronger assumption that $\mathcal{L} \geq 1$, 
  the functor $\CC^{\Gr} \to \CC^{\Fil}$ which is right adjoint to the associated graded functor
  restricts to an equivalence of symmetric monoidal $1$-categories
  \[ \CC^{\Gr, \mathcal{L}, \heartsuit} \simeq \CC^{\Fil, \mathcal{L}, \heartsuit}. \]
\end{lem}

\begin{ntn}
  We write $- \otimes^{\heartsuit, \mathcal{L}} -$ for the tensor product induced on either $\CC^{\Gr, \mathcal{L}, \heartsuit}$ or $\CC^{\Fil, \mathcal{L}, \heartsuit}$. In the case that $\mathcal{L} = \o$, we simply write $- \otimes ^{\heartsuit} -$.
\end{ntn}

%The $\mathcal{L}$-twisted $t$-structure is easy to understand in the graded case:
%
%If we assume that $\mathcal{L} \geq 0$, then we may obtain similar results in the filtered case.

\begin{proof}
  The statements (1), (2) and (3) are all clear. We therefore restrct ourselves to the filtered case for the rest of the proof.

  From the expression of the tensor product as a Day convolution, we observe that a tensor product of connective objects has $n^{\mathrm{th}}$ term presented as a colimit over a diagram of connective objects tensored with $\mathcal{L}$ at least $n$ times. Since $\mathcal{L} \geq 0$ we may conclude that (1') holds.

  The expression for the homotopy objects in (3') follows from (2') in a straightforward way.
  We now prove (2').
  Suppose that $X_\bullet$ is an object such that $\mathcal{L}^{\otimes -n} \otimes X_n < 0$.
  We want to show that $X_\bullet < 0$, i.e. that it recieves no maps from $Y_\bullet \geq 0$.
  Using the condition that $\mathcal{L} \geq 0$, we learn that $[Y_n, X_m] = 0$ for all $n > m$, which is enough to imply that $[Y_\bullet, X_\bullet] = 0$.

  Suppose that $X_\bullet < 0$; then we would like to show that $\mathcal{L}^{\otimes -n} \otimes X_n < 0$.
  Associated to every $Y \geq 0$ in $\CC$ we can consider the filtered object $Y \otimes \mathcal{L}^{\otimes n}(n)$, which is depicted below:
  \[ \dots \to 0 \to 0 \to Y \otimes \mathcal{L} ^{\otimes n} \xrightarrow{\id} Y \otimes \mathcal{L} ^{\otimes n} \xrightarrow{\id} \dots \]
Here, the first nonzero object occurs at position $n$.
  Using the asssumption that $\mathcal{L} \geq 0$, we can conclude that $Y \otimes \mathcal{L}^{\otimes n}(n) \geq 0$. Now, on mapping out of this object we have
  \[ * \simeq \Omega^{\infty} \Map_{\CC^{\Fil}} ( Y \otimes \mathcal{L}^{\otimes n}(n), X) \simeq \Omega^\infty \Map_\CC ( Y \otimes \mathcal{L}^{\otimes n}, X_n), \]
  which implies that $\mathcal{L}^{\otimes -n} \otimes X_n < 0$ for all $n$, as desired.

  We now proceed to prove the final statement.
  Let $\tau : \o_{\CC^{\Fil}} (-1) \to \o_{\CC^{\Fil}}$ denote the shift map in $\CC^{\Fil}$. Then $C\tau$ admits an $\E_\infty$-algebra structure in $\CC^{\Fil}$ such that there is an equivalence of symmetric monoidal categories $\CC^{\Gr} \simeq \Mod(\CC^{\Fil}; C\tau)$ (cf. \Cref{prop:mod-fil}). Using the assumption that $\mathcal{L} \geq 0$, we learn that $\o_{\CC^{\Fil}}$ and $C\tau$ are $\geq 0$. As a consequence, we find that $\CC^{\Gr, \mathcal{L}, \heartsuit}$ may be identified with $\Mod(\CC^{\Fil, \mathcal{L}, \heartsuit}; \piu_0 ^{\mathcal{L}} C\tau)$.

  The proposition will therefore follow if we prove that $\piu_{0} ^{\mathcal{L}} \o_{\CC^{\Fil}} \to \piu_0 ^{\mathcal{L}} C\tau$ is an equivalence.
  By \Cref{prop:twist-filt}, we find that
  \[(\piu_{0} ^{\mathcal{L}, \Fil} \o_{\CC^{\Fil}})_n \cong \begin{cases} \mathcal{L}^{\otimes n} \pi_0 (\mathcal{L} ^{\otimes -n}) & \text{ if } n \leq 0, \\ 0 & \text{ otherwise.} \end{cases}\]
    Since $\mathcal{L} \geq 1$ by assumption, we find for $n < 0$ that $\mathcal{L}^{\otimes n} \pi_0 \mathcal{L} ^{\otimes -n} \cong 0$, whereas for $n=0$ we just get $\pi_0 \o_{\CC}$.

  On the other hand, we have that
  \[(\piu_0 ^{\mathcal{L}, \Fil} C\tau)_n \cong \begin{cases} \pi_0 \o_{\CC} & \text{ if } n = 0, \\ 0 & \text{ otherwise.} \end{cases}\]
  It follows from this that the map $\piu_{0} ^{\mathcal{L}, \Fil} \o_{\CC^{\Fil}} \to \piu_0 ^{\mathcal{L}, \Fil} C\tau$ is an equivalence, as desired.
\end{proof}

We can now profitably define the twist functors.

\begin{cnstr} \label{cnstr:twist}
  Given a symmetric monoidal Grothendieck abelian $1$-category $\mathcal{A}$ and an invertible object $a \in \mathcal{A}$, we can construct a monoidal functor
  $i_a : \Z \to \mathcal{A}$ which sends $1$ to $a$.\footnote{By symmetric monoidal Grothendieck abelian $1$-category, we mean a commutative algebra object of the symmetric monoidal category $\mathrm{Groth}_1$ discussed in \Cref{rec:groth-ab}. In other words, we assume that the tensor product commutes with colimits in each variable.}
  If the swap map $s_{a,a}$ is the identity, then we can make this functor symmetric monoidal.
  Tensoring up with $\Abh$, we obtain a monoidal (or symmetric monoidal) functor
  \[ i_a : \Ab^{\Gr, \heartsuit} \xrightarrow{\o(1) \mapsto a} \mathcal{A}. \]
  
  In our situation of interest, we take $\mathcal{A} = \CC^{\Gr, \mathcal{L}, \heartsuit}$ and $a = (\piu_0 \o) \otimes \mathcal{L}(1)$. Then we can define the functor $\twist^{\mathcal{L}}$ as the composite
  \[ \CC^{\Gr, \heartsuit} \cong \CC^\heartsuit \otimes \Ab^{\Gr, \heartsuit} \xrightarrow{\mathrm{Id} \otimes i_{(\piu_0 \o) \otimes \mathcal{L}(1)}} \CC^\heartsuit \otimes \CC^{\Gr, \mathcal{L}, \heartsuit} \xrightarrow{\otimes} \CC^{\Gr, \mathcal{L}, \heartsuit}. \]
  On objects this has the effect of sending $\{X_n\}$ to $\{X_n \otimes \mathcal{L}^{\otimes n}\}$.
  It is easy to see that this is a monoidal equivalence between $\CC^{\Gr, \heartsuit}$ and $\CC^{\Gr, \mathcal{L}, \heartsuit}$, and is further symmetric monoidal if $s_{a, a}$ is the identity. 
\end{cnstr}

The above construction succinctly explains how much the symmetric monoidal structure on $\CC^{\Gr, \mathcal{L}, \heartsuit}$ has been twisted in terms of the quantity $s_{a, a}$. We now give a description of this quantity in terms of the following standard definition:

\begin{dfn}
  Given a dualizable object $X$,
  with dual $X^\vee$, unit $\eta$, and counit $\epsilon$,
  the trace of an endomorphism $f: X \to X$ is the element $\mathrm{tr}(f) \in \pi_0 \End(\o)$ given by the composite
  \[ \o \xrightarrow{\eta} X \otimes X^\vee \xrightarrow{ f \otimes X^\vee} X \otimes X^\vee \xrightarrow{s_{X,X^\vee}} X^{\vee} \otimes X \xrightarrow{\epsilon} \o. \]
  The trace of the identity is called the Euler characteristic and denoted $\chi(X)$.
\end{dfn}

A simple diagram chase tells us that the Euler characteristic is multiplicative and defines an $\End(\o)$-linear map $\mathrm{tr} : \End(X) \to \End(\o)$. Moreover, one can compute that the trace of the swap map $s_{X,X}$ is equal to $\chi(X)$. In the case where $X$ is invertible the trace map is an isomorphism, and we can use this to conclude that the condition $s_{a,a} = \mathrm{Id}_{a \otimes a}$ in \Cref{cnstr:twist} is equivalent to asking that $\chi(a) = 1$. Specializing further to the case $a = (\piu_0 \o) \otimes \mathcal{L}(1)$, we note that the diagram computing the Euler characteristic of $(\piu_0 \o) \otimes \mathcal{L}(1)$ is $\piu_0^{\mathcal{L},\Gr}$ applied to the diagram computing the Euler characteristic of $\mathcal{L}(1)$ in the ambient category $\CC^{\Gr}$. This implies that
\[ \chi(a) = \chi(\mathcal{L}(1)) = \chi(\mathcal{L})\chi(\o(1)) = \chi(\mathcal{L}) \cdot 1. \]
Thus, we have proved:

\begin{lem} \label{prop:twist}
  The functor $\twist^{\mathcal{L}}$ can be made symmetric monoidal if $\chi(\mathcal{L}) = 1$.
\end{lem}

%% \begin{lem}
%%   The assignment $\{X_n\} \mapsto \{X_n \otimes \mathcal{L}^{\otimes n}\}$ determines an automorphism of $\twist^{\mathcal{L}}$ of $\CC^{\Gr}$ which is $t$-exact if we equip the source with the standard $t$-structure and the target with the $\mathcal{L}$-twisted $t$-structure. 
%% \end{lem}

%% \begin{proof} Clear from \Cref{dfn:L-twist} and \Cref{prop:twist-grade}. \end{proof}

%% The question we would like to answer is ``when can these functors be equipped with a symmetric monoidal structure?''. The first obstruction to this comes from the Euler characteristic, which we define below.

Conversely, if we assume that $\twist^{\mathcal{L}}$ is symmetric monoidal, then we can make the following Euler characteristic computation:
\[ 1 = \twist^{\mathcal{L}} (1) = \twist^{\mathcal{L}}( \chi (\o(1))) = \chi ( \twist^{\mathcal{L}}(\o(1))) = \chi ( \mathcal{L}(1) ) = \chi (\mathcal{L}). \]
%% In fact, we can show that at the level of the heart this is the only obstruction to making the functor $\twist^{\mathcal{L}}$ symmetric monoidal. 

%% \begin{lem} \label{prop:twist}
%%   Since $\twist^{\mathcal{L}}$ is $t$-exact it naturally restricts to a functor,
%%   \[ \twist^{\mathcal{L}} : \CC^{\Gr, \heartsuit} \simeq \CC^{\Gr, \mathcal{L}, \heartsuit}. \]
%%   This functor may naturally be given the structure of a monoidal functor.
%%   Moreover, it may be given a symmetric monoidal structure if $\chi(\mathcal{L}) = 1$.
%% \end{lem}

%% \begin{proof}
%%    To give $\twist^{\mathcal{L}}$ the structure of a monoidal functor, we must specify natural isomorphisms
%%    \[\{X_n \otimes \mathcal{L}^{\otimes n}\} \otimes^{\heartsuit, \mathcal{L}} \{Y_n \otimes \mathcal{L}^{\otimes n}\} \cong \{(X_\bullet \otimes^{\heartsuit} Y_\bullet)_n \otimes \mathcal{L}^{\otimes n}\},\]
%%    i.e. specify natural isomorphisms
%%    \[\piu_0 ^{\mathcal{L}} (\{X_n \otimes \mathcal{L}^{\otimes n}\} \otimes \{Y_n \otimes \mathcal{L}^{\otimes n}\}) \cong \{(\piu_{0} (X_\bullet \otimes Y_\bullet))_n \otimes \mathcal{L}^{\otimes n}\}.\]
%%    Now, by \Cref{prop:twist-grade}, we have a natural isomorphism
%%    \[\{(\piu_{0} (X_\bullet \otimes Y_\bullet))_n \otimes \mathcal{L}^{\otimes n}\} \cong \piu_0 ^{\mathcal{L}} (\{(X_\bullet \otimes Y_\bullet)_n \otimes \mathcal{L}^{\otimes n}\}).\]
%%    It therefore suffices to specify natural isomorphisms
%%    \[\{X_n \otimes \mathcal{L}^{\otimes n}\} \otimes \{Y_n \otimes \mathcal{L}^{\otimes n}\} \cong \{(X_\bullet \otimes Y_\bullet)_n \otimes \mathcal{L}^{\otimes n}\}\]
%%    in the homotopy $1$-category $h\CC$ of $\CC$. From this point on we work in the symmetric monoidal $1$-category $h\CC$.

%%    Expanding out the definitions, we find that we need to provide natural isomorphisms
%%    \[\{\bigoplus_{a+b=n} (X_a \otimes \mathcal{L}^{\otimes a}) \otimes (Y_b \otimes \mathcal{L}^{\otimes b})\} \cong \{(\bigoplus_{a+b=n} X_a \otimes Y_b) \otimes \mathcal{L}^{\otimes n}\}.\]
%%    We now provide, for each $a$ and $b$, an equivalence
%%    \[(X_a \otimes \mathcal{L}^{\otimes a}) \otimes (Y_b \otimes \mathcal{L}^{\otimes b}) \cong (X_a \otimes Y_b) \otimes \mathcal{L}^{\otimes a+b},\]
%%    namely the swap map $s_{\mathcal{L}^{\otimes a}, Y_b}$ (here and throughout the rest of the proof we will suppress the associativity isomorphisms).

%%    It is not hard to check that this equips $\twist^{\mathcal{L}}$ with the structure of a monoidal functor.
%%    To check that it is in fact a symmetric monoidal functor, we need to check that the following diagram commutes for each $a$ and $b$:
%%    \begin{center}
%%      \begin{tikzcd} [column sep = 1.5in, row sep = large]
%%        (X_a \otimes \mathcal{L} ^{\otimes a}) \otimes (Y_b \otimes \mathcal{L}^{\otimes b}) \ar[r,"s_{(X_a \otimes \mathcal{L} ^{\otimes a}), (Y_b \otimes \mathcal{L}^{\otimes b})}"] \ar[d,"s_{\mathcal{L}^{\otimes a}, Y_b}"] & (Y_b \otimes \mathcal{L} ^{\otimes b}) \otimes (X_a \otimes \mathcal{L}^{\otimes a}) \ar[d, "s_{\mathcal{L}^{\otimes b}, X_a}"] \\
%%        (X_a \otimes Y_b) \otimes \mathcal{L}^{\otimes a+b} \ar[r,"s_{X_a, Y_b}"] & (Y_b \otimes X_a) \otimes \mathcal{L}^{\otimes a+b}.
%%      \end{tikzcd}
%%    \end{center}
%%    Tracing around the the bottom of the diagram, we obtain the map
%%    \[s_{X_a, \mathcal{L}^{\otimes b}} \circ s_{X_a, Y_b} \circ s_{\mathcal{L}^{\otimes a}, Y_b}.\]
%%    On the other hand, going across the top, we have
%%    \begin{align*}
%%      &s_{(X_a \otimes \mathcal{L} ^{\otimes a}), (Y_b \otimes \mathcal{L}^{\otimes b})} \\
%%      &= s_{(X_a \otimes \mathcal{L} ^{\otimes a}), \mathcal{L}^{\otimes b}} \circ s_{(X_a \otimes \mathcal{L} ^{\otimes a}), Y_b} \\
%%      &=  s_{X_a, \mathcal{L}^{\otimes b}} \circ s_{\mathcal{L} ^{\otimes a}, \mathcal{L}^{\otimes b}} \circ s_{X_a, Y_b} \circ s_{\mathcal{L} ^{\otimes a}, Y_b}.
%%    \end{align*}
%%    Now, we have $s_{\mathcal{L}^{\otimes a}, \mathcal{L}^{\otimes b}} = (s_{\mathcal{L}, \mathcal{L}})^{\circ (ab)}$ and therefore is will suffice to show that $s_{\mathcal{L}, \mathcal{L}}$ is the identity.

%%    The trace defines a $\pi_0\o$-linear isomorphism between $\pi_0End(\mathcal{L}^{\otimes 2})$ and $\pi_0\o$ which sends the identity to $\chi(\mathcal{L}^{\otimes 2}) = \chi( \mathcal{L} )^2$. One can then directly compute that the trace of $s_{\mathcal{L}, \mathcal{L}}$ is equal to $\chi ( \mathcal{L} )$. Now, using the assumption that $\chi (\mathcal{L}) = 1$ we may conclude.
%% \end{proof}

\begin{rmk}
  Examining the failure of $\twist^{\mathcal{L}}$ to be symmetric monoidal, we find that in general the symmetric monoidal structure on $\CC^{\Gr, \mathcal{L}, \heartsuit}$ is twisted by the Euler characteristic $\chi(\mathcal{L})$.
\end{rmk}
%
%\begin{rec}
%  Given a presentably symmetric monoidal category $\CC$, we may consider the Picard category $\Pic(\CC)$. This is the full subcategory
%
%  If $\CC$ is a $1$-category, then $\pi_n \pic (\CC) = 0$ for $n \neq 0,1$. Moreover, $\pi_0 \pic(\CC)$ is the Picard group of $\CC$ and $\pi_1 \Pic(\CC) \cong \aut(\o)$.
%\end{rec}
%
%\begin{prop}
%  The cofiber of the map $\pic (\CC^{\heartsuit}) \to \pic (\CC^{\Gr, \mathcal{L}, \heartsuit})$ is equivalent to $\Z$.
%\end{prop}
%
%\begin{proof}
%  Since $\CC^{\heartsuit} \to \CC^{\Gr, \mathcal{L}, \heartsuit}$ is fully faithful, the induced map  $\pic (\CC^{\heartsuit}) \to \pic (\CC^{\Gr, \mathcal{L}, \heartsuit})$
%\end{proof}
%
%\begin{prop}
%  The natural symmetric monoidal functor $\mathcal{P} (\Pic (\CC^{\Gr, \mathcal{L}, \heartsuit})) \otimes_{\mathcal{P} (\Pic (\CC^{\heartsuit}))} \CC^{\heartsuit} \to \CC^{\Gr, \mathcal{L}, \heartsuit}$ is an equivalence.
%\end{prop}
%
%As a consequence, the symmetric monoidal structure on $\CC^{\Gr, \mathcal{L}, \heartsuit}$ only depends on the map of spectra $\pic (\CC^{\heartsuit}) \to \pic (\CC^{\Gr, \mathcal{L}, \heartsuit})$.
%
%It follows that this map is determined by a map $\Z \to \Sigma \pic(\CC^{\heartsuit})$.
%\begin{lem}
%  Let $A$ and $B$ denote abelian groups, viewed as Eilenberg-Maclane spectra. Let $B[2] \subset B$ denote the subgroup of $2$-torsion elements. Then there are isomorphisms:
%  \[ [A,\Sigma B] \cong \Ext_{\Z} (A,B) \]
%  \[ [A, \Sigma^2 B] \cong \Hom_{\Z} (A,B[2]). \]
%\end{lem}
%
%We therefore need to understand the structure of $\pic(\CC^{\heartsuit})$.
%This has two nonzero homotopy groups: $\pi_0 \pic(\CC^{\heartsuit})$ is the Picard group of $\CC^{\heartsuit}$, and $\pi_1 \pic(\CC^{\heartsuit}) \cong \aut(\o_{\CC^{\heartsuit}})$.
%The single nonzero $k$-invariant $\pi_0 \pic(\CC^{\heartsuit}) \to \Sigma^{2} \pi_1 \pic(\CC^{\heartsuit}) \simeq \Sigma^{2} \aut(\o_{\CC^{\heartsuit}})$
%corresponds to the trace map $\Tr : \pi_0 \pic(\CC^{\heartsuit}) \to \aut(\o_{\CC^{\heartsuit}})[2]$ sending a Picard element $\mathcal{L}$ to its trace $\Tr(\mathcal{L})$.
%
%
%  The negative objects are easy to recognize in the graded case.
%  If we assume that $\mathcal{L}$ is $\geq 0$, then the negative are similarly easy to determine,
%  \begin{itemize}
%  \item An object $X_\bullet$ of $\mathcal{C}^{\Gr}$ is $\leq 0$ in the $\mathcal{L}$-twisted $t$-structure if and only if for all $n$, $\mathcal{L}^{\otimes -n} \otimes X_n$ is $\leq 0$ in the original $t$-structure.
%  \item Assuming that $\mathcal{L}$ is $\geq 0$, an object $X_\bullet$ of $\mathcal{C}^{\Fil}$ is $\leq 0$ in the $\mathcal{L}$-twisted $t$-structure if and only if its underlying graded object satisfies the same.
%  \end{itemize}
%
%In the graded case we have a simple formula for homotopy group functor,
%\[ (\underline{\pi}_0^{\mathcal{L}, \Gr} (X_\bullet) )_n \cong \mathcal{L}^{\otimes n} \underline{\pi}_0(\mathcal{L}^{\otimes -n} \otimes X_n) \]
%In the filtered case we have a similar formula as long as $\mathcal{L}$ is $\geq 0$,
%\[ (\underline{\pi}_0^{\mathcal{L}, \Fil} (X_\bullet) )_n \cong \mathcal{L}^{\otimes n} \underline{\pi}_0(\mathcal{L}^{\otimes -n} \otimes X_n) \]
%If we make the stronger assumption that $\mathcal{L}$ is $\geq 1$, then the shift map on an object of the heart becomes zero (up to contractible ambiguity) so $\underline{\pi}_0^{\mathcal{L}, \Fil}$ factors through the graded version of the heart with an equivalence
%\[ \underline{\pi}_0^{\mathcal{L},\Gr}(\Gr(-)) \cong \underline{\pi}_0^{\mathcal{L},\Fil}(-). \]
%Altogether, this lets us conclude that as long as $\mathcal{L}$ is $\geq 1$, there are equivalences
%\[ \mathcal{C}^{\Fil, \mathcal{L}\heartsuit} \cong \mathcal{C}^{\Gr, \mathcal{L}\heartsuit} \cong (\mathcal{C}^\heartsuit)^{\Gr}. \]
%where the second equivalence sends $X_\bullet$ to $\mathcal{L}^{\otimes \bullet} \otimes X_\bullet$.
%
%In the situation of \cite{dfn:L-twist} the compatibility of the $t$-structure with the tensor product turns $\mathcal{C}^\heartsuit$ into a symmetric monoidal $1$-category. Similarly, the assumption that $\mathcal{L}$ is $\geq 0$ ensures that the $\mathcal{L}$-twisted $t$-structure is compatible with the tensor product on filtered/graded objects. Using the equivalence of 1-categories from the previous paragraph, we have constructed a symmetric moniodal structure on $(\mathcal{C}^\heartsuit)^{\Gr}$ for each suitable $\mathcal{L}$. For later use we would like to study these symmetric monoidal structures.
%
%A monoidal $\mathcal{C}^\heartsuit$-linear functor with source $\mathcal{C}^{\heartsuit,\Gr}$ is equivalent to the choice of Picard element in the target (to which $\o(1)$ is sent). From this, we can conclude that the $\mathcal{L}$-twisted symmetric monoidal structure on $\mathcal{C}^{\heartsuit, \Gr}$ has the same underlying monoidal structure as the untwisted case.
%
%\begin{lem}
%  If we consider $\mathcal{C}^{\heartsuit, \Gr}$ as a monoidal $\mathcal{C}^\heartsuit$-algebra, then symmetric monoidal refinements of the monoidal structure are in bijective correspondence with 2-torsion elements of $\pi_0(\o)^\times$. The correspondence sends a symmetric monoidal refinement to $Tr(\o(1))$.
%\end{lem}
%
%\begin{proof}
%  \todo{write this.}
%\end{proof}
%

\begin{exm}
  If we take $\mathcal{C} = \Sp$ with its usual $t$-structure, then the $\Ss^1$-twisted category $\Sp^{\Gr, \Ss^1, \heartsuit}$ is equivalent to the symmetric monoidal category of graded abelian groups with the Koszul sign convention, since $\chi(\Ss^1) = -1$.

  On the other hand, $\Sp^{\Gr, \Ss^2, \heartsuit}$ is symmetric monoidally equivalent to the category of graded abelian groups by \Cref{prop:twist}, since $\chi(\Ss^2) = 1$. \todo{Define $\chi$ above, note that it is $2$-torsion for a Picard element.}
\end{exm}
%
%In [Ctau], the fact that $S^2$-twisted graded abelian groups are equivalent to graded abelian groups is implicitly used in proving their symmetric monoidal equivalence.
%\todo{Should we throw some shade about lazy non-proofs here?}

%
%Extending the example above, the case of interest for this paper is $\Ss^{\rho}$-twisted graded $\uZ$-module Mackey functors. Observe that since $[C_2] - 1$ maps to $1$ in $\uZ$ the twist is trivial (like the case over $\C$).

\begin{lem} \label{lem:D-twist}
  Suppose that $\CC$ is equivalent to the derived category of its heart $\CC^{\heartsuit}$ as a symmetric monoidal category with compatible $t$-structure. Then there exists an equivalence of monoidal categories $\twist^{\mathcal{L}} : \CC^{\Gr} \simeq \CC^{\Gr}$ making the following diagram of monoidal functors commute:
  \begin{center}
    \begin{tikzcd}
      \CC^{\Gr, \heartsuit} \ar[r, "\twist^{\mathcal{L}}"] \ar[d] & \CC^{\Gr, \mathcal{L}, \heartsuit} \ar[d] \\
      \CC^{\Gr} \ar[r, "\twist^{\mathcal{L}}"] & \CC^{\Gr}.
    \end{tikzcd}
  \end{center}
  If $\chi(\mathcal{L}) = 1$, then the above square naturally lifts to a diagram of symmetric monoidal functors.
  On objects, $\twist^{\mathcal{L}} : \CC^{\Gr} \to \CC^{\Gr}$ is given by $\{X_n\} \mapsto \{X_n \otimes \mathcal{L}^{\otimes n}\}$.
\end{lem}

\begin{proof}
  The assumption implies that $\CC^{\Gr}$, equipped with the usual $t$-structure, is equivalent as a symmetric monoidal category with compatible $t$-struture to $\mathcal{D}(\CC^{\Gr, \heartsuit})$.
  It also implies that $\CC^{\Gr}$, when equipped with the $\mathcal{L}$-twisted $t$-structure, is equivalent as a symmteric monoidal category with compatible $t$-struture to $\mathcal{D}(\CC^{\Gr, \mathcal{L}, \heartsuit})$.

  It follows that the equivalence of monoidal $1$-categories $\twist^{\mathcal{L}} : \CC^{\Gr, \heartsuit} \to \CC^{\Gr, \mathcal{L}, \heartsuit}$ of \Cref{prop:twist} determines a compatible equivalence of monoidal categories $\twist^{\mathcal{L}}: \CC^{\Gr} \to \CC^{\Gr}$, via taking derived categories.
  If $\chi(\mathcal{L}) = 1$, these are equivalences of symmetric monoidal categories by \Cref{prop:twist}.
\end{proof}

